\documentclass[a4paper,12pt]{report}

% =========================================================================
% GÓI HỖ TRỢ
% =========================================================================
\usepackage[utf8]{inputenc}
\usepackage[T5]{fontenc}
\usepackage[vietnamese]{babel}
\usepackage{times}

\usepackage{graphicx}
\usepackage{geometry}
\geometry{left=3cm, right=2cm, top=2cm, bottom=2cm}
\usepackage{float}
\usepackage{setspace}
\usepackage{booktabs}
\usepackage{longtable}
\usepackage{enumitem}
\usepackage{array}
\usepackage{xcolor}
\usepackage{listings}
\usepackage{caption}
\usepackage{subcaption}
\usepackage{tikz}
\usepackage{eso-pic}
\usepackage{amsmath, amssymb}
\usepackage{hyperref}
\hypersetup{
    colorlinks=true,
    linkcolor=blue,
    urlcolor=blue,
    citecolor=blue
}

% =========================================================================
% KHUNG VIỀN TRANG BÌA
% =========================================================================
\newcommand\BackgroundPic{%
\put(0,0){%
\parbox[b][\paperheight]{\paperwidth}{%
\vfill\centering
\begin{tikzpicture}
\draw[line width=3pt] (2cm,2cm) rectangle (\paperwidth-2cm,\paperheight-2cm);
\draw[line width=1pt] (2.15cm,2.15cm) rectangle (\paperwidth-2.15cm,\paperheight-2.15cm);
\end{tikzpicture}
\vfill}}}

% =========================================================================
% CẤU HÌNH LISTINGS (hiển thị code C)
% =========================================================================
\lstset{
    language=C,
    basicstyle=\ttfamily\small,
    keywordstyle=\color{blue}\bfseries,
    commentstyle=\color{gray}\itshape,
    stringstyle=\color{red!70!black},
    numbers=left,
    numberstyle=\tiny\color{gray},
    frame=single,
    breaklines=true,
    captionpos=b,
    tabsize=2
}

\begin{document}
\onehalfspacing

% =========================================================================
%                            TRANG BÌA
% =========================================================================
\begin{titlepage}
    \AddToShipoutPicture*{\BackgroundPic}
    \centering

    \vspace*{0.5cm}
    {\Large \textbf{ĐẠI HỌC BÁCH KHOA HÀ NỘI}} \\
    \vspace{0.2cm}
    {\large \textbf{TRƯỜNG CÔNG NGHỆ THÔNG TIN VÀ TRUYỀN THÔNG}} \\
    \vspace{1cm}

    \begin{figure}[h!]
        \centering
        \includegraphics[width=0.25\textwidth]{logo.png}
    \end{figure}

    \vspace{0.8cm}
    {\large \textbf{BÁO CÁO BÀI TẬP LỚN}} \\
    \vspace{0.3cm}
    {\large \textbf{HỆ NHÚNG --- IT4210}} \\ 
    \vspace{1.5cm}
    {\LARGE \textbf{ĐỀ TÀI:}} \\
    \vspace{0.5cm}
    {\Large \textbf{GAME ĐỐI KHÁNG 2D NARUTO VS BLEACH\\
    TRÊN VI ĐIỀU KHIỂN STM32F429ZIT6}} \\
    \vspace{1.2cm}

    \begin{table}[h]
        \centering
        \begin{tabular}{l l}
            \textbf{Giảng viên hướng dẫn:} & Đỗ Công Thuần \\ 
            & \\
            \textbf{Sinh viên thực hiện:}
                & Hoàng Đức Anh - [MSSV] \\  % <-- Điền MSSV
                & Nguyễn Việt Dũng - [MSSV] \\
                & Nguyễn Đức Hiếu - [MSSV] \\
                & Nguyễn Khánh Toàn - 20235847 \\
        \end{tabular}
    \end{table}

    \vfill
    {\large \textbf{Hà Nội, \today}}
    \vspace*{1cm}
\end{titlepage}

% =========================================================================
%                          LỜI CẢM ƠN
% =========================================================================
\chapter*{Lời cảm ơn}
\addcontentsline{toc}{chapter}{Lời cảm ơn}

Chúng em xin gửi lời cảm ơn chân thành đến thầy \textbf{[Tên giảng viên]} đã tận tình hướng dẫn, định hướng kỹ thuật và tạo điều kiện để nhóm được thực hiện một đề tài mang tính thực hành cao trong khuôn khổ môn học Hệ thống Nhúng.

Dự án này đòi hỏi nhóm phải làm chủ đồng thời nhiều lĩnh vực kỹ thuật: lập trình vi điều khiển bare-metal, giao tiếp ngoại vi tốc độ cao, xử lý đồ họa thời gian thực và thiết kế hệ thống trò chơi. Nhờ có sự định hướng từ thầy cùng với sự nỗ lực của các thành viên, nhóm đã hoàn thành được sản phẩm hoạt động đầy đủ trên phần cứng thực tế.

Chúng em cũng cảm ơn cộng đồng mã nguồn mở và các tài liệu kỹ thuật của STMicroelectronics đã cung cấp nền tảng kiến thức để nhóm nghiên cứu và triển khai.

Dù đã rất cố gắng, báo cáo chắc chắn vẫn còn những thiếu sót. Nhóm rất mong nhận được góp ý từ thầy để cải thiện trong các dự án tiếp theo.

\vspace{1cm}
\begin{flushright}
    \textit{Hà Nội, \today} \\
    \textbf{Nhóm sinh viên thực hiện}
\end{flushright}
\newpage

% =========================================================================
%                          TÓM TẮT
% =========================================================================
\chapter*{Tóm tắt}
\addcontentsline{toc}{chapter}{Tóm tắt}

Báo cáo này trình bày quá trình thiết kế và xây dựng một trò chơi đối kháng 2D theo phong cách \textit{Naruto vs Bleach} chạy trực tiếp trên vi điều khiển \textbf{STM32F429ZIT6} (ARM Cortex-M4, 16\,MHz) mà không sử dụng hệ điều hành nhúng (bare-metal).

Sản phẩm hỗ trợ chế độ chơi Người--Máy: người chơi điều khiển nhân vật qua joystick analog và bốn nút nhấn vật lý, trong khi đối thủ được điều phối bởi một bộ AI đơn giản chạy ngay trên MCU. Màn hình hiển thị 320$\times$240 pixel (ILI9341 tích hợp sẵn trên kit, giao tiếp FSMC 16-bit) với tốc độ làm mới logic trận đấu khoảng 30 khung/giây, sử dụng kỹ thuật \textit{dirty-rect} để tối ưu tốc độ vẽ.

Hệ thống gồm bốn nhân vật chiến đấu được hoạt hình hóa (Naruto, Sasuke, Ichigo, Vizard Ichigo) với đầy đủ cơ chế: di chuyển, nhảy, dash, combo tấn công ba đòn, đỡ đòn, kỹ năng đặc biệt, hitbox/hurtbox, tính toán sát thương và thanh HP. Toàn bộ dữ liệu sprite được mã hóa định dạng RGB565 với color-key trong suốt và lưu trực tiếp trong Flash của MCU.

\textbf{Từ khoá:} STM32F429, game đối kháng nhúng, ILI9341, sprite RGB565, dirty-rect, AI đối thủ, bare-metal C.

\newpage

% =========================================================================
%                       MỤC LỤC / DANH SÁCH HÌNH
% =========================================================================
\tableofcontents
\newpage
\listoffigures
\listoftables
\newpage

% =========================================================================
%               CHƯƠNG 1: GIỚI THIỆU ĐỀ TÀI
% =========================================================================
\chapter{Giới thiệu đề tài}

\section{Mô tả sản phẩm và mục tiêu}

Đề tài xây dựng một \textbf{trò chơi đối kháng 2D} lấy cảm hứng từ tựa game Flash nổi tiếng \textit{Bleach vs Naruto}, chạy hoàn toàn trên vi điều khiển \textbf{STM32F429ZIT6} (kit STM32F429I-DISCO) mà không sử dụng hệ điều hành nhúng hay GPU rời.

Mục tiêu kép của dự án gồm:
\begin{enumerate}[label=\arabic*)]
    \item \textbf{Mục tiêu kỹ thuật:} Chứng minh khả năng hiển thị đồ họa sprite thời gian thực, xử lý input đa kênh và thực thi logic game phức tạp (vật lý, va chạm, AI) hoàn toàn trên một MCU nhúng chi phí thấp.
    \item \textbf{Mục tiêu sản phẩm:} Tạo ra một trò chơi hoàn chỉnh, có thể chơi được trên phần cứng thực, với trải nghiệm người dùng tương đối mượt mà (không bị giật khung hình rõ rệt).
\end{enumerate}

Sản phẩm hỗ trợ chế độ \textbf{Người--Máy (PvAI)}: người chơi cầm joystick và bốn nút vật lý, đối thủ được điều khiển bởi một bộ AI phản ứng theo trạng thái trận đấu.

\section{Yêu cầu chức năng}

\subsection{Đầu vào}
\begin{itemize}
    \item \textbf{Joystick analog 2 trục} (PA0--PA1, ADC1): điều khiển hướng di chuyển (trái/phải) và đỡ đòn (kéo xuống).
    \item \textbf{Bốn nút nhấn vật lý} (PB0--PB3, GPIOB, active-low, pull-up nội):
    \begin{itemize}
        \item PB0 --- Tấn công thường (Attack)
        \item PB1 --- Nhảy (Jump)
        \item PB2 --- Kỹ năng đặc biệt (Skill)
        \item PB3 --- Lướt nhanh (Dash)
    \end{itemize}
\end{itemize}

\subsection{Đầu ra}
\begin{itemize}
    \item \textbf{Màn hình LCD ILI9341} 320$\times$240 pixel, màu RGB565, giao tiếp SPI.
    \item Hiển thị: nền arena 2D, hai nhân vật sprite hoạt ảnh, thanh HP của cả hai bên, tên trạng thái hành động hiện tại.
\end{itemize}

\subsection{Chức năng game}

\begin{table}[H]
\centering
\caption{Danh sách chức năng của game}
\begin{tabular}{c p{4cm} p{8cm}}
\toprule
\textbf{STT} & \textbf{Chức năng} & \textbf{Mô tả} \\
\midrule
1 & Di chuyển & Chạy trái/phải, dừng về Idle tự động \\
2 & Nhảy & Nhảy có trọng lực, di chuyển ngang trên không \\
3 & Lướt (Dash) & Lướt nhanh theo hướng đang mặt \\
4 & Đỡ đòn (Block) & Giảm sát thương và knockback khi bị tấn công từ phía trước \\
5 & Tấn công combo & Chuỗi 3 đòn liên tiếp (attack1 $\to$ attack2 $\to$ attack3), kích hoạt bằng cách nhấn nút trong cửa sổ combo \\
6 & Kỹ năng đặc biệt & Mỗi nhân vật có 1 kỹ năng riêng (ví dụ: Getsuga Tensho cho Vizard Ichigo, Chidori cho Sasuke) \\
7 & Hitbox / Hurtbox & Mỗi trạng thái có vùng va chạm gây và nhận sát thương riêng biệt \\
8 & Thanh HP & HP bắt đầu 100, giảm khi trúng đòn; khi HP = 0, nhân vật ngã và trận đấu kết thúc \\
9 & AI đối thủ & CPU phân tích khoảng cách, trạng thái, HP để ra quyết định (tiếp cận, tấn công, đỡ, dùng kỹ năng) \\
10 & Màn hình chọn nhân vật & Lưới 6 nhân vật, hiển thị avatar và banner nhân vật được chọn \\
\bottomrule
\end{tabular}
\end{table}

\subsection{Nhân vật}

Game có \textbf{4 nhân vật có thể chiến đấu}, mỗi nhân vật có bộ hoạt ảnh và hitbox riêng:

\begin{table}[H]
\centering
\caption{Danh sách nhân vật chiến đấu}
\begin{tabular}{c l l}
\toprule
\textbf{STT} & \textbf{Nhân vật} & \textbf{Kỹ năng đặc biệt} \\
\midrule
1 & Naruto & Rasengan (đòn cận chiến tầm gần, mạnh) \\
2 & Sasuke & Chidori (đâm tới trước, có hiệu ứng điện) \\
3 & Ichigo & Getsuga Tensho (phóng projectile) \\
4 & Vizard Ichigo & Getsuga Tensho + Dịch chuyển tức thời (teleport) \\
\bottomrule
\end{tabular}
\end{table}

\section{Yêu cầu phi chức năng}

\begin{table}[H]
\centering
\caption{Yêu cầu phi chức năng}
\begin{tabular}{c p{3.5cm} p{8.5cm}}
\toprule
\textbf{STT} & \textbf{Tiêu chí} & \textbf{Mô tả / Ngưỡng đạt} \\
\midrule
1 & Tốc độ khung hình & Logic game cập nhật mỗi 33\,ms ($\approx$30 FPS); màn hình không bị giật rõ rệt trong điều kiện bình thường \\
2 & Độ trễ đầu vào & Phản hồi nút nhấn trong vòng 1 chu kỳ game tick (33\,ms) \\
3 & Độ ổn định & Không xảy ra crash (HardFault, stack overflow) trong thời gian chơi liên tục ít nhất 10 phút \\
4 & Chứa dữ liệu trong Flash & Toàn bộ sprite, ảnh nền, UI asset phải vừa trong Flash 2\,MB của STM32F429ZIT6 \\
5 & Kiến trúc bare-metal & Hệ thống được thiết kế theo mô hình super-loop có tick cố định, không phụ thuộc RTOS. Đây là lựa chọn chủ động: game có luồng xử lý tuần tự đơn giản, không cần lập lịch đa nhiệm; bare-metal giúp kiểm soát timing chính xác hơn và tránh overhead context switch của RTOS \\
6 & Độ hoàn thiện UI & Màn hình splash, menu chọn nhân vật, màn hình game over phải hiển thị đúng và đầy đủ \\
7 & Khả năng mở rộng & Kiến trúc phần mềm cho phép thêm nhân vật mới chỉ bằng cách bổ sung file \texttt{*\_moveset.c/h} mà không sửa engine \\
\bottomrule
\end{tabular}
\end{table}

% =========================================================================
%               CHƯƠNG 2: THIẾT KẾ HỆ THỐNG
% =========================================================================
\chapter{Thiết kế hệ thống}

\section{Phân chia chức năng phần cứng -- phần mềm}

Game đối kháng được xây dựng theo kiến trúc \textbf{bare-metal} (không RTOS), trong đó toàn bộ logic chạy trong một vòng lặp chính \texttt{while(1)} với chu kỳ tick cố định 33\,ms. Bảng~\ref{tab:hw-sw} phân chia rõ trách nhiệm giữa phần cứng và phần mềm.

\begin{table}[H]
\centering
\caption{Phân chia chức năng phần cứng -- phần mềm}
\label{tab:hw-sw}
\begin{tabular}{p{4cm} p{4.5cm} p{5.5cm}}
\toprule
\textbf{Chức năng} & \textbf{Phần cứng thực hiện} & \textbf{Phần mềm thực hiện} \\
\midrule
Đọc joystick & ADC1 (12-bit, PA0--PA1) & Hiệu chỉnh điểm trung tâm, deadzone, ánh xạ sang lệnh game \\
Đọc nút nhấn & GPIO GPIOB (PB0--PB3), pull-up & Phát hiện cạnh lên (edge detect), tổng hợp bitmask input \\
Hiển thị hình ảnh & LCD ILI9341 (FSMC, 320$\times$240, onboard kit) & Điều khiển FSMC, thuật toán dirty-rect, vẽ sprite RGB565 \\
Định thời game & SysTick (HAL\_GetTick) & Kiểm soát chu kỳ tick 33\,ms, đảm bảo logic không trôi lệch \\
Logic trận đấu & --- (hoàn toàn SW) & Vật lý, va chạm hitbox/hurtbox, sát thương, AI đối thủ \\
Lưu trữ sprite & Flash 2\,MB (STM32F429ZIT6) & Script Python ngoại tuyến mã hoá sprite sang mảng C/RGB565 \\
\bottomrule
\end{tabular}
\end{table}

\section{Thiết kế phần cứng}

\subsection{Danh sách linh kiện và module}

\begin{table}[H]
\centering
\caption{Danh sách linh kiện chính}
\label{tab:components}
\begin{tabular}{c p{3cm} p{6cm} p{3.5cm}}
\toprule
\textbf{STT} & \textbf{Linh kiện} & \textbf{Thông số chính} & \textbf{Vai trò} \\
\midrule
1 & STM32F429ZIT6 & ARM Cortex-M4, 2\,MB Flash, 256\,KB RAM, LQFP144 & Vi điều khiển trung tâm \\
2 & LCD ILI9341 & 320$\times$240 px, RGB565, FSMC 16-bit (tích hợp kit) & Hiển thị game \\
3 & Joystick analog & 2 trục, đầu ra $0\sim3.3$\,V & Điều hướng nhân vật \\
4 & Nút nhấn & Active-low, 4 nút vật lý & Tấn công / Nhảy / Kỹ năng / Dash \\
\bottomrule
\end{tabular}
\end{table}


\subsection{Sơ đồ ghép nối}

Hình~\ref{fig:hw-block} mô tả sơ đồ khối phần cứng tổng thể.

\begin{figure}[H]
\centering
\begin{tikzpicture}[
    box/.style={draw, rounded corners, minimum width=3cm, minimum height=1.2cm,
                text centered, font=\small, align=center},
    arr/.style={->, thick, >=stealth}
]

% ---- Khung kit STM32F429I-DISCO (bao quanh MCU + LCD) ----
\node[draw, dashed, rounded corners, fill=blue!5,
      minimum width=8.5cm, minimum height=4cm,
      label={[font=\footnotesize\itshape]north:STM32F429I-DISCO Kit}]
      (kit) at (2, 0) {};

% ---- MCU bên trong kit ----
\node[box, fill=blue!20, minimum width=3.8cm, minimum height=2.5cm]
      (mcu) at (0, 0) {
        \textbf{STM32F429ZIT6}\\
        \footnotesize Cortex-M4 @ 16\,MHz\\
        \footnotesize 2\,MB Flash / 256\,KB RAM
      };

% ---- LCD bên trong kit ----
\node[box, fill=orange!15, minimum width=3cm, minimum height=2cm]
      (lcd) at (4.2, 0) {
        \textbf{LCD ILI9341}\\
        \footnotesize 320$\times$240, RGB565\\
        \footnotesize (tích hợp kit)
      };

% ---- Mũi tên nội bộ MCU -> LCD (FSMC) ----
\draw[arr] (mcu.east) -- node[above, font=\tiny]{FSMC 16-bit} (lcd.west);

% ---- Ngoại vi ngoài kit ----
\node[box, fill=green!15] (joy) at (-4.5, 1.3)
      {Joystick\\[-2pt]\footnotesize analog 2 trục};
\node[box, fill=green!15] (btn) at (-4.5,-1.3)
      {4 Nút nhấn\\[-2pt]\footnotesize vật lý};
\node[box, fill=gray!20, minimum width=3cm] (pwr) at (2,-3.2)
      {Nguồn 3.3\,V\\[-2pt]\footnotesize USB / ST-LINK};

% ---- Mũi tên ngoài vào kit ----
\draw[arr] (joy.east) -- node[above, font=\tiny]{ADC1 (PA0, PA1)} (mcu.west |- joy.east);
\draw[arr] (btn.east) -- node[above, font=\tiny]{GPIO (PB0--PB3)} (mcu.west |- btn.east);
\draw[arr] (pwr.north) -- (kit.south -| pwr.north);

\end{tikzpicture}
\caption{Sơ đồ ghép nối phần cứng hệ thống}
\label{fig:hw-block}
\end{figure}


\subsection{Cấu hình chân và tốc độ giao tiếp}

\begin{table}[H]
\centering
\caption{Cấu hình chân MCU}
\label{tab:pins}
\begin{tabular}{c p{2.2cm} p{2.5cm} p{3cm} p{4cm}}
\toprule
\textbf{STT} & \textbf{Chân MCU} & \textbf{Ngoại vi} & \textbf{Chế độ} & \textbf{Ghi chú} \\
\midrule
1 & PA0 & ADC1 CH0 & Analog in & Trục X joystick \\
2 & PA1 & ADC1 CH1 & Analog in & Trục Y joystick \\
3 & PB0 & GPIOB & Digital in, pull-up & Nút Attack \\
4 & PB1 & GPIOB & Digital in, pull-up & Nút Jump \\
5 & PB2 & GPIOB & Digital in, pull-up & Nút Skill \\
6 & PB3 & GPIOB & Digital in, pull-up & Nút Dash \\
7+ & FSMC Bank 1 & ILI9341 (onboard) & FSMC 16-bit parallel & Tích hợp sẵn trên kit, driver xử lý tự động \\
\bottomrule
\end{tabular}
\end{table}

\noindent\textbf{Tốc độ giao tiếp FSMC:} LCD ILI9341 được kết nối qua FSMC Bank~1 của STM32F429I-DISCO theo giao thức Intel 8080 16-bit. Ở xung HCLK = 16\,MHz, thông lượng lý thuyết đạt $\approx 32$\,MB/s, thời gian ghi một frame đầy đủ 320$\times$240$\times$2\,byte $\approx 4{,}8$\,ms. Tuy nhiên do overhead câu lệnh và độ trễ FSMC, trên thực tế thời gian vẽ toàn màn hình vẫn đủ lớn để thuật toán dirty-rect mang lại lợi ích đáng kể về hiệu suất.

\noindent\textbf{ADC:} Cấu hình đơn kênh, sampling time 480 cycle, phân giải 12-bit ($0 \sim 4095$). Điểm trung tâm joystick hiệu chỉnh khi khởi động bằng trung bình 16 mẫu.

\section{Thiết kế phần mềm}

\subsection{Sơ đồ khối phần mềm}

\begin{figure}[H]
\centering
\begin{tikzpicture}[
    layer/.style={draw, minimum width=11cm, minimum height=0.85cm,
                  text centered, font=\small},
    arr/.style={->, thick, >=stealth}
]
\node[layer, fill=red!15]    (L1) at (0, 5.0)
    {\textbf{Application:} \texttt{battle\_demo.c} --- Vòng lặp game, AI, điều phối rendering};
\node[layer, fill=orange!20] (L2) at (0, 3.8)
    {\textbf{Game Logic:} \texttt{combat\_actor.c}, \texttt{combat\_attack\_data.c}, \texttt{combat\_input.c}};
\node[layer, fill=yellow!30] (L3) at (0, 2.6)
    {\textbf{Asset:} \texttt{*\_moveset.c}, \texttt{chidori\_data.c}, \texttt{choose\_ui\_assets.c}};
\node[layer, fill=green!20]  (L4) at (0, 1.4)
    {\textbf{Graphics:} \texttt{sprite\_render.c}, \texttt{lcd\_port.c}, \texttt{game\_ui.c}};
\node[layer, fill=blue!15]   (L5) at (0, 0.2)
    {\textbf{Driver:} \texttt{ILI9341\_STM32\_Driver.c}, \texttt{ILI9341\_GFX.c}, STM32 HAL};
\node[layer, fill=gray!25]   (L6) at (0,-1.0)
    {\textbf{Hardware:} STM32F429ZIT6, ILI9341 LCD (onboard), Joystick, Buttons};

\draw[arr] (L1.south) -- (L2.north);
\draw[arr] (L2.south) -- (L3.north);
\draw[arr] (L3.south) -- (L4.north);
\draw[arr] (L4.south) -- (L5.north);
\draw[arr] (L5.south) -- (L6.north);
\end{tikzpicture}
\caption{Sơ đồ khối phần mềm theo lớp}
\label{fig:sw-block}
\end{figure}


\subsection{Biểu đồ luồng hoạt động hệ thống}

\begin{figure}[H]
\centering
\begin{tikzpicture}[
    proc/.style={draw, rectangle, rounded corners, minimum width=5cm, minimum height=0.75cm, text centered, font=\small},
    dec/.style={draw, diamond, aspect=3, text centered, font=\small, inner sep=1pt},
    arr/.style={->, thick},
    startstop/.style={draw, ellipse, minimum width=3cm, minimum height=0.65cm, font=\small\bfseries, fill=black!80, text=white}
]
\node[startstop] (start) at (0, 0)    {Khởi động MCU};
\node[proc, fill=blue!10]   (init)  at (0,-1.2) {HAL\_Init() + SystemClock\_Config()};
\node[proc, fill=blue!10]   (ginit) at (0,-2.3) {BattleDemo\_Init(): LCD, ADC, GPIO, AI, Background cache};
\node[proc, fill=orange!10] (upd)   at (0,-3.4) {BattleDemo\_Update()};
\node[dec,  fill=yellow!30] (chk)   at (0,-4.7) {Đủ 33\,ms?};
\node[proc, fill=green!10]  (inp)   at (0,-6.0) {Đọc Input (ADC joystick + GPIO buttons)};
\node[proc, fill=green!10]  (ai)    at (0,-7.1) {AI: Observe() $\to$ ChooseAction() $\to$ bitmask input CPU};
\node[proc, fill=green!10]  (phy)   at (0,-8.2) {CombatActor\_Update(): physics, state machine, frame index};
\node[proc, fill=green!10]  (col)   at (0,-9.3) {Kiểm tra va chạm Hitbox $\cap$ Hurtbox, áp dụng sát thương};
\node[proc, fill=green!10]  (prj)   at (0,-10.4){Cập nhật Projectile (Getsuga, Chidori)};
\node[proc, fill=red!10]    (drw)   at (0,-11.5){Dirty-Rect render $\to$ LCD\_Port\_Flush()};

\draw[arr] (start)--(init); \draw[arr] (init)--(ginit);
\draw[arr] (ginit)--(upd);  \draw[arr] (upd)--(chk);
\draw[arr] (chk.east) -- node[right,font=\tiny]{Không} ++(2.5,0) |- (upd.east);
\draw[arr] (chk.south) -- node[right,font=\tiny]{Có} (inp.north);
\draw[arr] (inp)--(ai); \draw[arr] (ai)--(phy);
\draw[arr] (phy)--(col); \draw[arr] (col)--(prj); \draw[arr] (prj)--(drw);
\draw[arr] (drw.south) -- ++(0,-0.3) -- ++(-5,0) |- (upd.west);
\end{tikzpicture}
\caption{Biểu đồ luồng hoạt động của hệ thống game}
\label{fig:flow}
\end{figure}

\subsection{Mô tả các module phần mềm chính}

\subsubsection{Module đầu vào (\texttt{combat\_input})}
Đọc ADC1 kênh CH0/CH1 (joystick) và GPIO GPIOB (nút nhấn), trả về \texttt{uint8\_t} bitmask. Joystick áp dụng deadzone $\pm650$\,LSB; nút nhấn chỉ kích hoạt trên \textbf{cạnh lên} để tránh lặp liên tục.

\subsubsection{Module quản lý nhân vật (\texttt{combat\_actor})}
Mỗi nhân vật là một \texttt{CombatActor} gồm vị trí, vận tốc, HP, trạng thái. Máy trạng thái 9 nút (\texttt{IDLE, RUN, JUMP, DASH, ATTACK, SKILL, BLOCK, HIT, DEAD}) với vật lý đơn giản: trọng lực $g=1$\,px/tick, lực nhảy $v_y=-11$, tốc độ rơi tối đa $9$\,px/tick, biên arena $x\in[32,288]$.

\subsubsection{Module dữ liệu chiến đấu (\texttt{combat\_attack\_data})}
Lưu bảng hurtbox (theo trạng thái) và hitbox (theo frame, attackStep) cho từng nhân vật. Va chạm kiểm tra bằng phép giao AABB:
\[
\text{Overlap}(A,B) \;\Leftrightarrow\; A.x < B.x{+}B.w \;\wedge\; A.x{+}A.w > B.x \;\wedge\; A.y < B.y{+}B.h \;\wedge\; A.y{+}A.h > B.y
\]

\subsubsection{Module hoạt ảnh (\texttt{*\_moveset.c})}
Mỗi nhân vật có file độc lập chứa pixel RGB565 (color-key \texttt{0xF81F} = trong suốt), pivot tâm chân, thời gian frame. Dữ liệu được sinh tự động bằng script Python từ spritesheet nguồn.

\subsubsection{Module vẽ Dirty-Rect (\texttt{sprite\_render} + \texttt{battle\_demo})}
Chỉ vẽ lại vùng hình chữ nhật bao frame cũ và frame mới. Mỗi pixel trong vùng đó: nếu là pixel sprite hợp lệ thì vẽ sprite, ngược lại lấy từ bộ nhớ đệm nền (\texttt{s\_backgroundPixels[]}) đã được tạo khi khởi động, tránh đọc lại Flash.

\subsubsection{Module AI (\texttt{BattleAiController})}
Hoạt động theo chu kỳ hành động với 2 bước:
\begin{itemize}
    \item \textbf{Quan sát:} Cập nhật trạng thái đối thủ vào buffer \texttt{pending}; sau mỗi 220\,ms mới chuyển sang \texttt{visible} (mô phỏng độ trễ phản xạ người thật).
    \item \textbf{Quyết định:} Tại 65\% chu kỳ hành động, AI chọn lệnh tiếp theo theo thứ tự ưu tiên: Phản đòn $>$ Đỡ $>$ Kỹ năng $>$ Tấn công $>$ Tiếp cận/Lùi, có thêm yếu tố ngẫu nhiên (LCG RNG) để tránh máy móc.
\end{itemize}


% =========================================================================
%          CHƯƠNG 3: CÀI ĐẶT VÀ XÂY DỰNG HỆ THỐNG
% =========================================================================
\chapter{Cài đặt và xây dựng hệ thống}

\section{Mã nguồn và hướng dẫn cài đặt}

Toàn bộ mã nguồn của dự án được lưu trữ công khai trên GitHub:

\begin{center}
    \url{https://github.com/hoangducanh2005/bleach_vs_naruto_STM32F429ZI}
\end{center}

\noindent File \texttt{README.md} tại thư mục gốc mô tả đầy đủ các bước cài đặt và yêu cầu công cụ. Tóm tắt:

\begin{table}[H]
\centering
\caption{Công cụ phát triển}
\begin{tabular}{l l l}
\toprule
\textbf{Công cụ} & \textbf{Phiên bản} & \textbf{Mục đích} \\
\midrule
STM32CubeIDE & 1.17.0 & Biên dịch và nạp firmware \\
STM32CubeMX & 6.17.0 & Cấu hình ngoại vi MCU \\
STM32Cube FW\_F4 & V1.28.3 & HAL driver STM32F4 \\
Python & 3.10+ & Chạy script chuyển đổi sprite \\
Pillow (PIL) & 10.x & Xử lý ảnh trong script Python \\
Git & 2.x & Quản lý phiên bản mã nguồn \\
\bottomrule
\end{tabular}
\end{table}

\noindent\textbf{Các bước cài đặt:}
\begin{enumerate}
    \item Clone repository: \texttt{git clone https://github.com/hoangducanh2005/bleach\_vs\_naruto\_STM32F429ZI}
    \item Mở project trong STM32CubeIDE (File $\to$ Import $\to$ Existing Projects).
    \item Kết nối kit STM32F429I-DISCO qua cổng USB (ST-LINK).
    \item Nhấn \textbf{Build} (\texttt{Ctrl+B}) rồi \textbf{Run} để nạp firmware và chạy game.
    \item (Tùy chọn) Chạy các script trong thư mục \texttt{tools/} để tái sinh dữ liệu sprite nếu cần thay đổi tài nguyên đồ họa.
\end{enumerate}

\section{Mô tả các module phần mềm chính}

\subsection{Cấu trúc thư mục}

\begin{table}[H]
\centering
\caption{Cấu trúc thư mục dự án}
\begin{tabular}{p{4.5cm} p{9cm}}
\toprule
\textbf{Thư mục / File} & \textbf{Mô tả} \\
\midrule
\texttt{Core/Src/} & Toàn bộ mã nguồn C của game và driver \\
\texttt{Core/Inc/} & File header tương ứng \\
\texttt{assets/} & Sprite nguồn (PNG), ảnh giao diện chọn tướng \\
\texttt{assets/raw/} & Spritesheet gốc trước khi xử lý \\
\texttt{tools/} & Script Python chuyển đổi sprite sang mảng C \\
\texttt{docs/} & Tài liệu thiết kế và báo cáo \\
\bottomrule
\end{tabular}
\end{table}

\subsection{Các module cốt lõi}

\begin{table}[H]
\centering
\caption{Danh sách module phần mềm cốt lõi}
\begin{tabular}{p{4.5cm} p{9cm}}
\toprule
\textbf{Module (file .c/.h)} & \textbf{Chức năng chính} \\
\midrule
\texttt{battle\_demo} & Vòng lặp game chính: điều phối input, AI, physics, rendering, HUD \\
\texttt{combat\_actor} & Máy trạng thái nhân vật, vật lý, xử lý combo, trạng thái stun \\
\texttt{combat\_input} & Đọc joystick (ADC) và nút nhấn (GPIO), trả về bitmask \\
\texttt{combat\_attack\_data} & Bảng hitbox/hurtbox theo từng trạng thái và frame của nhân vật \\
\texttt{sprite\_render} & Thuật toán dirty-rect, composite sprite lên nền, gọi LCD driver \\
\texttt{*\_moveset} & Dữ liệu hoạt ảnh RGB565 của từng nhân vật (Naruto, Sasuke, Ichigo, Vizard) \\
\texttt{game\_ui} & Vẽ màn hình chọn nhân vật, main menu, HUD, màn hình kết thúc \\
\texttt{ILI9341\_STM32\_Driver} & Driver LCD cấp thấp: khởi tạo ILI9341, ghi pixel qua FMC \\
\texttt{lcd\_port} & Lớp trừu tượng LCD (adapter pattern), cách ly driver khỏi game \\
\bottomrule
\end{tabular}
\end{table}

\section{Đóng góp của từng thành viên}

\begin{table}[H]
\centering
\caption{Phân công và đóng góp của từng thành viên}
\begin{tabular}{p{3.5cm} p{3.5cm} p{7cm}}
\toprule
\textbf{Thành viên} & \textbf{Số commit (nếu có)} & \textbf{Nội dung đóng góp chính} \\
\midrule
Hoàng Đức Anh & [Điền số commit] & [Điền đóng góp chính tại đây] \\[4pt]
Nguyễn Việt Dũng & [Điền số commit] & [Điền đóng góp chính tại đây] \\[4pt]
Nguyễn Đức Hiếu & [Điền số commit] & [Điền đóng góp chính tại đây] \\[4pt]
Nguyễn Khánh Toàn & [Điền số commit] & [Điền đóng góp chính tại đây] \\
\bottomrule
\end{tabular}
\end{table}

\noindent\textit{Lưu ý: Số commit không phản ánh hoàn toàn mức độ đóng góp thực tế của từng thành viên do có nhiều phần công việc nghiên cứu và lập trình nhóm thực hiện trực tiếp cùng nhau hoặc kiểm thử trên phần cứng thực tế mà không tạo commit riêng lẻ.}

\section{Kết quả}

\subsection{Kết quả đạt được}

Sản phẩm đã được xây dựng hoàn chỉnh và chạy ổn định trên phần cứng thực (kit STM32F429I-DISCO). Các tính năng đã hoàn thành:

\begin{itemize}
    \item Màn hình splash, menu chính, chọn nhân vật (6 nhân vật), chọn độ khó.
    \item Trận đấu Người--Máy hoàn chỉnh với 4 nhân vật chiến đấu được hoạt ảnh hoá.
    \item Cơ chế combat đầy đủ: combo 3 đòn, kỹ năng đặc biệt (Chidori, Getsuga Tensho), projectile.
    \item AI đối thủ phản ứng theo trạng thái trận đấu với 3 mức độ khó.
    \item Màn hình Game Over và kết quả trận đấu.
\end{itemize}

\subsection{Video và ảnh demo}

\begin{itemize}
    \item \textbf{Video demo:} \url{[Điền link video demo tại đây]}  % <-- Điền link
    \item \textbf{Ảnh chụp màn hình:} Xem thư mục \texttt{docs/} trong repository.
\end{itemize}

\section{Đánh giá}

\subsection{Đánh giá theo yêu cầu đề ra}

\begin{table}[H]
\centering
\caption{Đánh giá mức độ hoàn thành yêu cầu}
\begin{tabular}{c p{5cm} c p{5cm}}
\toprule
\textbf{STT} & \textbf{Yêu cầu} & \textbf{Đạt?} & \textbf{Ghi chú} \\
\midrule
1 & Game chạy được trên STM32F429 & \checkmark & Ổn định, không crash \\
2 & Hiển thị sprite hoạt ảnh & \checkmark & 4 nhân vật, đầy đủ trạng thái \\
3 & Điều khiển joystick + nút & \checkmark & Latency $\leq$ 33\,ms \\
4 & AI đối thủ & \checkmark & Phản ứng theo khoảng cách, HP \\
5 & Kỹ năng đặc biệt & \checkmark & Chidori, Getsuga, teleport \\
6 & UI đầy đủ & \checkmark & Splash, select, game over \\
7 & Dữ liệu vừa Flash 2\,MB & \checkmark & Tổng $\approx 1.8$\,MB \\
8 & Không dùng RTOS & \checkmark & Bare-metal super-loop \\
\bottomrule
\end{tabular}
\end{table}

\subsection{Ưu điểm}

\begin{itemize}
    \item \textbf{Kiến trúc module hóa cao:} Mỗi nhân vật độc lập trong một file \texttt{*\_moveset.c}, có thể thêm nhân vật mới mà không sửa engine.
    \item \textbf{Hiệu suất tốt với bare-metal:} Thuật toán dirty-rect giúp tốc độ khung hình đạt 30 FPS logic trên MCU chỉ 16\,MHz.
    \item \textbf{Quy trình asset tự động:} Script Python tự động trích xuất và biên dịch sprite, giảm sai sót thủ công.
    \item \textbf{AI có độ trễ phản xạ:} Mô phỏng phản ứng người thật, tránh cảm giác AI máy móc.
\end{itemize}

\subsection{Nhược điểm và hướng cải thiện}

\begin{itemize}
    \item \textbf{Chưa hỗ trợ 2 người chơi:} Hiện chỉ có chế độ Người--Máy; cần thêm joystick thứ hai cho PvP.
    \item \textbf{Flash gần đầy:} Với $\approx 1.8$\,MB dữ liệu sprite, khả năng thêm nhân vật mới bị giới hạn bởi dung lượng Flash 2\,MB.
    \item \textbf{AI đơn giản:} AI hiện tại chưa học từ hành vi người chơi; có thể cải thiện bằng thuật toán heuristic phức tạp hơn.
    \item \textbf{Không có âm thanh:} Phiên bản hiện tại chưa tích hợp buzzer hay DAC để phát âm thanh hiệu ứng.
\end{itemize}

% =========================================================================
%          CHƯƠNG 4: AI-ASSISTED DISCLAIMER
% =========================================================================
\chapter{Tuyên bố sử dụng công cụ AI}

\section{Khẳng định sử dụng AI}

Nhóm \textbf{xác nhận [CÓ / KHÔNG] sử dụng} các công cụ trí tuệ nhân tạo (AI) trong quá trình thực hiện project và viết báo cáo này.

\section{Danh sách công cụ AI đã sử dụng}

\begin{table}[H]
\centering
\caption{Công cụ AI được sử dụng trong dự án}
\begin{tabular}{l l p{7cm}}
\toprule
\textbf{Công cụ} & \textbf{Nhà cung cấp} & \textbf{Mục đích sử dụng} \\
\midrule
Công cụ AI 1 & [Nhà cung cấp 1] & [Mô tả chi tiết mục đích sử dụng 1] \\
Công cụ AI 2 & [Nhà cung cấp 2] & [Mô tả chi tiết mục đích sử dụng 2] \\
Công cụ AI 3 & [Nhà cung cấp 3] & [Mô tả chi tiết mục đích sử dụng 3] \\
\bottomrule
\end{tabular}
\end{table}

\section{Mô tả chi tiết kỹ thuật prompt}

Các thành viên điền chi tiết các prompt tiêu biểu đã sử dụng vào các mục dưới đây để minh họa:

\subsection{Prompt hỗ trợ lập trình và sửa lỗi (Coding/Debugging)}

\begin{itemize}
    \item \textbf{Thành viên thực hiện:} [Tên thành viên]
    \item \textbf{Mục đích:} [Mô tả mục đích sử dụng prompt]
    \item \textbf{Nội dung prompt:}
    \begin{quote}
    \textit{``[Điền nội dung prompt viết bằng tiếng Anh/Việt vào đây]''}
    \end{quote}
    \item \textbf{Kết quả phản hồi của AI:} [Mô tả ngắn gọn kết quả phản hồi hoặc giải pháp AI cung cấp]
\end{itemize}

\subsection{Prompt hỗ trợ thiết kế hệ thống và giải thuật}

\begin{itemize}
    \item \textbf{Thành viên thực hiện:} [Tên thành viên]
    \item \textbf{Mục đích:} [Mô tả mục đích sử dụng prompt]
    \item \textbf{Nội dung prompt:}
    \begin{quote}
    \textit{``[Điền nội dung prompt viết bằng tiếng Anh/Việt vào đây]''}
    \end{quote}
    \item \textbf{Kết quả phản hồi của AI:} [Mô tả ngắn gọn kết quả phản hồi hoặc giải pháp AI cung cấp]
\end{itemize}

\subsection{Prompt hỗ trợ viết tài liệu và báo cáo}

\begin{itemize}
    \item \textbf{Thành viên thực hiện:} [Tên thành viên]
    \item \textbf{Mục đích:} [Mô tả mục đích sử dụng prompt]
    \item \textbf{Nội dung prompt:}
    \begin{quote}
    \textit{``[Điền nội dung prompt viết bằng tiếng Anh/Việt vào đây]''}
    \end{quote}
    \item \textbf{Kết quả phản hồi của AI:} [Mô tả ngắn gọn kết quả phản hồi hoặc giải pháp AI cung cấp]
\end{itemize}

\section{Mức độ và giới hạn sử dụng AI}

\begin{itemize}
    \item \textbf{Phần tự làm (Không sử dụng AI):} [Điền các nội dung nhóm tự thực hiện hoàn toàn, ví dụ: tự viết logic C bare-metal, tự test phần cứng, thiết kế hitbox, cấu hình ngoại vi trên CubeMX...]
    \item \textbf{Phần AI hỗ trợ:} [Điền các phần có sự trợ giúp của AI, ví dụ: dịch tài liệu, định dạng bảng biểu và sơ đồ TikZ trong LaTeX, tra cứu cú pháp...]
    \item \textbf{Giới hạn chịu trách nhiệm:} Nhóm cam kết đã kiểm tra tính chính xác của mọi thông tin do AI cung cấp và chịu trách nhiệm hoàn toàn về mặt kỹ thuật đối với sản phẩm thực tế.
\end{itemize}

% =========================================================================
% HẾT BÁO CÁO
% =========================================================================

\end{document}
